\documentclass[11pt]{article}
\usepackage[margin=1in]{geometry}
\usepackage{amsmath,amssymb}
\usepackage{graphicx}
\usepackage{booktabs}
\usepackage{siunitx}
\usepackage{array}
\usepackage{url}
\usepackage[hidelinks]{hyperref}
\usepackage{xcolor}
\graphicspath{{figures/chaos/}}

\title{Phase-Space Volume Contraction, Not the 0--1 Test, Separates\\
Conservative from Dissipative Chaos: A Cross-System Cautionary Result}
\author{Melissa Ellison\\\small Independent Researcher}
\date{\small 2026-06-19}

\begin{document}
\maketitle

\begin{abstract}
Deterministic chaos defeats long-horizon trajectory prediction, but many structural
features of a chaotic system remain measurable. We test whether a scalar diagnostic
can distinguish conservative Hamiltonian chaos from dissipative chaos, rather than
merely detecting chaos itself. Across a 26-system benchmark, an earlier version of
this study appeared to show that the Gottwald--Melbourne 0--1 statistic separates
conservative from dissipative systems. A trajectory-length audit falsified that
conclusion: for chaotic Hamiltonian systems the statistic is length-dependent and
can climb from near zero to near one (double pendulum: $K\approx-0.13$ at 4k steps
to $K\approx0.99$ at 16k steps, mean over 4 initial conditions). The 0--1 statistic
is therefore a chaos detector, not a conservative/dissipative classifier. The robust
discriminator is instead the phase-space volume-contraction rate
$\Sigma\lambda=\langle\nabla\!\cdot f\rangle$, equal to the sum of Lyapunov
exponents: zero to within $10^{-3}$ for all 8 Hamiltonian systems and strictly
negative for all 11 finite-dimensional dissipative flows, matching analytic traces
where available. We show why the false 0--1 split is seductive: the entire
$K$-vs-length curve is predicted from the observable autocovariance by a closed-form
Green--Kubo expression (Pearson $r=0.955$ against measured $K$). A held-out
symbolic-itinerary experiment further shows that attractor transition grammar is
predictable ($0.770$ vs.\ $0.233$ baseline on cell changes, 21/23 chaotic systems
beating baseline) even when microstate forecasts are Lyapunov-limited. The
contribution is a corrected discriminator, a reproducible warning about finite-time
misuse of a widely used chaos statistic \cite{falconer2007}, and a bounded
constructive result about predicting the climate, not the weather.
\end{abstract}

\section{Introduction}

Chaotic systems share sensitivity to initial conditions, yet they do not share the
same geometric or physical structure. A Hamiltonian double pendulum preserves
phase-space volume and remains on an energy shell; the Lorenz system
\cite{lorenz1963} contracts onto a strange attractor. A scalar chaos detector can report that both systems are
chaotic without answering whether the chaos is conservative or dissipative.

This paper asks a narrow question: among low-cost diagnostics---those requiring only
a scalar time series, not the full Lyapunov spectrum or the model vector
field---what separates conservative Hamiltonian chaos from dissipative chaos? We
tested two candidates. The first was the Gottwald--Melbourne 0--1 test
\cite{gottwald2004,gottwald2009,falconer2007}, which maps a scalar time series to a
statistic $K$ approaching one for chaotic dynamics and zero for regular dynamics.
The second was phase-space volume contraction,
$\Sigma\lambda=\langle\nabla\!\cdot f\rangle$, the divergence of the vector field
averaged along the trajectory, equal to the sum of Lyapunov exponents for smooth
flows \cite{benettin1980,kantz2004}. Liouville's theorem makes the latter a
natural discriminator: Hamiltonian flows preserve phase-space volume, while
dissipative flows contract it \cite{strogatz1994}.

The central result is a negative one. Early runs at short trajectory lengths
suggested that conservative chaos had $K\approx0$ while dissipative chaos had
$K\approx1$. Longer trajectories and initial-condition ensembles overturned that
reading. We therefore present the correction as part of the result: the 0--1 test
remains useful as a chaos detector, but it should not be used as a
conservative/dissipative classifier without additional structure.

\section{Systems and measurements}

The benchmark contains 26 systems spanning maps, dissipative flows, Hamiltonian
flows, delay dynamics, and non-chaotic controls (Table~\ref{tab:systems}). Systems
include the Lorenz~63 flow \cite{lorenz1963}, the H\'{e}non--Heiles potential
\cite{henon1964}, and the Mackey--Glass delay equation \cite{mackey1977}, among
others listed in Table~\ref{tab:systems}. Flows are integrated by fixed-step RK4
after warm-up onto the attractor; maps are iterated directly. Mechanical Hamiltonian
systems are integrated in canonical $(q,p)$ coordinates. Known analytic invariants
are tracked as integration hygiene checks.

\begin{table}[ht]
\centering
\small
\caption{The 26-system benchmark. $\lambda_1$: largest Lyapunov exponent (positive
= chaotic); ``n.c.''\ means not computed via global algorithm (Mackey--Glass
requires delay-embedding; known value from the literature is $\lambda_1\approx0.007$
for $\tau=17$). $K$: 0--1 statistic from the primary initial condition at 4000 steps
($n_{\mathrm{steps,flow}}$); for slow-mixing Hamiltonians the ensemble mean over
four initial conditions can differ substantially (double pendulum: $K=-0.13$
at 4000 steps, mean over four ICs; see Fig.~\ref{fig:klen}). $\Sigma\lambda$: measured
phase-space volume contraction (N/A for maps and the delay system, which lack a
finite-dimensional smooth vector field).}
\label{tab:systems}
\begin{tabular}{llr r r r}
\toprule
System & Type & Dim & $\lambda_1$ & $K$ & $\Sigma\lambda$ \\
\midrule
\multicolumn{6}{l}{\textit{Dissipative flows}} \\
Lorenz 63          & Dissipative & 3  & $+0.761$ & $0.996$ & $-13.667$ \\
R\"{o}ssler        & Dissipative & 3  & $+0.057$ & $0.801$ & $-5.737$  \\
Duffing (forced)   & Dissipative & 3  & $+0.114$ & $0.729$ & $-0.300$  \\
Lorenz 96 ($N=5$)  & Dissipative & 5  & $+0.587$ & $0.944$ & $-5.000$  \\
Aizawa             & Dissipative & 3  & $+0.122$ & $0.927$ & $-0.183$  \\
Sprott B           & Dissipative & 3  & $+0.217$ & $0.999$ & $-1.000$  \\
Halvorsen          & Dissipative & 3  & $+0.856$ & $0.954$ & $-4.200$  \\
Thomas             & Dissipative & 3  & $+0.046$ & $0.995$ & $-0.625$  \\
Chen               & Dissipative & 3  & $+2.134$ & $0.998$ & $-10.000$ \\
Ueda (forced)      & Dissipative & 3  & $+0.100$ & $0.513$ & $-0.050$  \\
Lorenz 96 ($N=10$) & Dissipative & 10 & $+1.129$ & $0.996$ & $-10.000$ \\
\midrule
\multicolumn{6}{l}{\textit{Hamiltonian flows}} \\
Double pendulum          & Hamiltonian & 4 & $+0.888$ & $0.770$ & $0.000$ \\
H\'{e}non--Heiles        & Hamiltonian & 4 & $+0.026$ & $0.877$ & $0.000$ \\
CR3BP (Earth--Moon)      & Hamiltonian & 4 & $+0.168$ & $-0.046$ & $0.000$ \\
Yang--Mills ($x^2y^2$)   & Hamiltonian & 4 & $+0.241$ & $-0.110$ & $0.000$ \\
Pullen--Edmonds          & Hamiltonian & 4 & $+0.317$ & $-0.100$ & $0.000$ \\
Quartic coupled          & Hamiltonian & 4 & $+0.250$ & $0.242$  & $0.000$ \\
Spring pendulum          & Hamiltonian & 4 & $+0.116$ & $0.979$  & $0.000$ \\
Coupled Duffing (Ham.)   & Hamiltonian & 4 & $+0.214$ & $0.957$  & $0.000$ \\
\midrule
\multicolumn{6}{l}{\textit{Maps (no vector field; $\Sigma\lambda$ not applicable)}} \\
Logistic ($r=3.9$)       & Map & 1 & $+1.322$ & $1.000$ & N/A \\
H\'{e}non map            & Map & 2 & $+0.584$ & $0.999$ & N/A \\
Standard map ($\kappa=1.2$) & Map & 2 & $+1.652$ & $0.996$ & N/A \\
\midrule
\multicolumn{6}{l}{\textit{Delay system ($\Sigma\lambda$ not applicable)}} \\
Mackey--Glass ($\tau=17$) & Delay & $\infty$ & n.c.$^{\dagger}$ & $0.855$ & N/A \\
\midrule
\multicolumn{6}{l}{\textit{Non-chaotic controls}} \\
Harmonic oscillator      & Control & 2 & $0.000$ & $0.599$ & N/A \\
Quasiperiodic 2-torus    & Control & 2 & $0.000$ & $0.961$ & N/A \\
Random walk              & Control & 2 & ---      & $0.045$ & N/A \\
\bottomrule
\end{tabular}
\end{table}

\paragraph{0--1 test.}
For a scalar observable $\phi(j)$ and frequency $c$, define translation variables
\begin{equation}
p_c(n)=\sum_{j=1}^{n}\phi(j)\cos(jc),\qquad
q_c(n)=\sum_{j=1}^{n}\phi(j)\sin(jc).
\end{equation}
The mean-square displacement over lag $n$ is
\begin{equation}
M_c(n)=\frac{1}{N}\sum_{j=1}^{N}\bigl[(p_c(j+n)-p_c(j))^2+(q_c(j+n)-q_c(j))^2\bigr].
\label{eq:Mc_def}
\end{equation}
We use $K_c=\mathrm{corr}(\log n,\log M_c(n))$, averaged over $n_c=12$ random
frequencies $c$ uniformly drawn from $(0,\pi)$.  This log--log form is equivalent
to the linear-regression coefficient used by Gottwald and Melbourne~\cite{gottwald2009}
in the asymptotic regime where $M_c(n)\sim n$ for chaotic and $M_c(n)\sim\mathrm{const}$
for regular dynamics, but it can assign high $K$ to quasiperiodic orbits near resonance
(where $M_c(n)\sim n^2$); this point is noted in Section~\ref{sec:disc}.
In the asymptotic regime, $K\to1$ for chaotic dynamics and $K\to0$ for regular
dynamics. The scalar observable is the first coordinate of the state vector for
flows; for maps the natural scalar output is used. The finite-length regime is the
failure mode studied here.

\paragraph{Phase-space volume contraction.}
For a vector field $\dot{x}=f(x)$,
\begin{equation}
\Sigma\lambda = \left\langle \nabla\!\cdot f(x(t))\right\rangle
=\left\langle\sum_i \frac{\partial f_i}{\partial x_i}\right\rangle.
\end{equation}
For smooth flows this equals the sum of Lyapunov exponents. Hamiltonian flows have
$\nabla\!\cdot f=0$ in canonical coordinates; dissipative flows contract volume.
Where analytic traces are available, they are used as cross-checks against the
numerical central-difference estimator.

\section{Results}

\subsection{The 0--1 statistic is not a conservative/dissipative discriminator}

On sufficiently long trajectories, Hamiltonian $K$ values scatter across the same
range as dissipative chaos. The decisive evidence is the length audit
(Fig.~\ref{fig:klen}). Dissipative systems such as Lorenz and Chen read
$K\approx1$ at all tested lengths. Hamiltonian systems can begin near zero and
climb toward one as the trajectory length grows. The double pendulum moves from
$K\approx-0.13$ (mean over 4 initial conditions) at 4k steps to $K\approx0.99$ at
16k steps; the coupled Duffing Hamiltonian moves from $K\approx-0.14$ at 4k to
$K\approx0.94$ at 16k. This is the pattern one
expects from a chaos detector whose finite-time behavior is governed by observable
mixing time, not by whether the system is conservative
(Section~\ref{sec:gk}).

\begin{figure}[h]
\centering
\includegraphics[width=0.78\textwidth]{fig2_K_vs_length}
\caption{The 0--1 statistic is trajectory-length dependent for slow-mixing
Hamiltonian systems. Purple curves can climb from near zero to near one; cyan
dissipative systems read near one at all lengths. Dotted curves are the
closed-form Green--Kubo predictions.}
\label{fig:klen}
\end{figure}

\subsection{Volume contraction separates the classes}

The phase-space contraction rate cleanly separates the finite-dimensional systems
with defined vector fields (Fig.~\ref{fig:disc}). The three maps and the
Mackey--Glass delay system are excluded from this analysis because they lack a
finite-dimensional smooth vector field; $\Sigma\lambda$ is defined only for ordinary
differential equation systems. Among those, all 8 Hamiltonian flows have
$\Sigma\lambda=0$ to within $10^{-3}$, and all 11 dissipative flows have
$\Sigma\lambda<0$. The estimator reproduces analytic traces where they
exist: Lorenz63 gives $-13.667=-\sigma-1-\beta$, Chen gives $-10$, Lorenz96$_N$
gives $-N$, Sprott~B gives $-1$, Thomas gives $-3b$, and Ueda gives $-k$. Because
$\Sigma\lambda$ is a vector-field property rather than a finite scalar-series
random walk, it has no trajectory-length crossover analogous to the 0--1 statistic.

\begin{figure}[h]
\centering
\includegraphics[width=0.72\textwidth]{fig1_sigma_lambda_discriminator}
\caption{Phase-space volume contraction separates conservative from dissipative
chaos. Hamiltonian systems sit at zero; dissipative flows contract. The horizontal
axis is symlog so both weak and strong contraction are visible.}
\label{fig:disc}
\end{figure}

The Hamiltonian cases also pass integration hygiene: all 8 systems conserve
an analytic invariant (energy, or the Jacobi integral for CR3BP), and the relative
drift in that invariant is at most $7\times10^{-4}$ over the run. The measured
largest Lyapunov exponent is positive for every chaotic Hamiltonian system. Thus the
zero volume contraction is not a failure to reach chaos; it is the expected
Liouville structure.

\subsection{Why the 0--1 failure happens}
\label{sec:gk}

The finite-time behavior of $K$ is not mysterious. For a mean-zero stationary
observable $\phi$ (mean-subtracted in the implementation) with autocovariance
$C(k)=E[\phi(j)\phi(j+k)]$, expanding Eq.~\eqref{eq:Mc_def} and applying
stationarity yields the exact mean-square displacement
\begin{equation}
M_c(n)=\sum_{k=-(n-1)}^{n-1}(n-|k|)C(k)\cos(kc)
      =nC(0)+2\sum_{k=1}^{n-1}(n-k)C(k)\cos(kc).
\label{eq:mc}
\end{equation}
Its asymptotic diffusion coefficient is the Green--Kubo / Wiener--Khinchin power
spectral density at frequency $c$ \cite{green1954,kubo1957}:
\begin{equation}
D_c=\lim_{n\to\infty}\frac{M_c(n)}{n}
=C(0)+2\sum_{k\ge1}C(k)\cos(kc).
\label{eq:dc}
\end{equation}
Since $K_c$ is the log--log correlation computed from $M_c(n)$, the
observable autocovariance predicts the full length-dependent curve. In the
artifacts, the predicted $K$ from Eq.~\eqref{eq:mc} matches measured $K$ with
Pearson $r=0.955$ and mean absolute error $0.077$ (Fig.~\ref{fig:gk}); the
predicted crossover length is within one ladder step on 20/23 chaotic systems.
The apparent class split was therefore a mixing-time artifact.

\begin{figure}[h]
\centering
\includegraphics[width=0.52\textwidth]{fig3_green_kubo_scatter}
\caption{Closed-form Green--Kubo prediction of measured 0--1 $K$ from the
observable autocovariance alone.}
\label{fig:gk}
\end{figure}

\subsection{Held-out itinerary grammar: predict the climate, not the weather}

The same observatory also measures symbolic attractor grammar. Each attractor is
coarsened into $k=6$ cells by k-means on the first half of the trajectory
(results are robust to $k\in\{4,6,8\}$; all values are in the data artifacts); an
order-1 Markov model is then evaluated on the held-out second half. Raw next-symbol
accuracy across 23 chaotic systems is $0.925$, but most of that is trivial
dwell-time persistence: finely sampled flows tend to remain in the same cell, and
the persistence baseline is $0.863$. The non-trivial result conditions on the
roughly 14\% of steps where the cell changes. On those transitions, the learned
grammar predicts the destination cell with mean accuracy $0.770$ versus a $0.233$
baseline, beating baseline on 21/23 systems (Fig.~\ref{fig:itin}); the remaining
2 (Duffing forced, Ueda forced) had zero cell transitions in the held-out half at
$k=6$ and are excluded from the transition mean. This is not
microstate prediction. It is a bounded, held-out statement that chaotic attractors
can have learnable symbolic route grammar.

\begin{figure}[h]
\centering
\includegraphics[width=0.74\textwidth]{fig5_itinerary_prediction}
\caption{Held-out next-cell prediction when the cell changes. The transition
grammar beats the most-common-transition baseline on 21/23 chaotic systems.}
\label{fig:itin}
\end{figure}

\subsection{Growth-rate diagnostics: M$_c$ slope analysis}
\label{sec:slope}

The log--log correlation $K_c$ collapses the entire $M_c(n)$ trajectory to a
single scalar. A complementary view is the \emph{growth-rate exponent}
$\alpha = \mathrm{d}\log M_c/\mathrm{d}\log n$, estimated by OLS on the
c-averaged $\log M_c(n)$ over $N_{\rm lag}=30$ geometrically-spaced lags
(data: \path{results/chaos_slope_analysis.json}).
All trajectories use $N=12{,}000$ integration steps---three times the 4,000-step
primary runs of Table~\ref{tab:systems}---to expose asymptotic behaviour, with
the same decimation procedure as Section~\ref{sec:gk}. The long-range sub-slope
$\alpha_{\mathrm{lo}}$ ($n \ge N_{\mathrm{dec}}/2$) additionally characterises
whether growth saturates at longer lags.

\paragraph{$\alpha$--$K$ decoupling for Hamiltonian systems.}
Four of eight Hamiltonian systems show a clear decoupling: CR3BP ($\alpha=0.847$),
Pullen--Edmonds ($\alpha=0.694$), quartic coupled ($\alpha=0.888$), and spring
pendulum ($\alpha=0.972$) all reach the diffusive regime while recording
$K\in[-0.17,\,-0.07]$ in the same run (Table~\ref{tab:slope}, bold rows).
The most striking example is the spring pendulum: starting from the identical
initial condition used in Table~\ref{tab:systems}, $K$ falls from $0.979$ at
4,000 steps to $-0.096$ at 12,000 steps, while $\alpha=0.972$ confirms
$M_c$ is growing diffusively throughout both runs. The mechanism is direct.
$\alpha$ is estimated from the mean $M_c(n)$ averaged over all 12 frequencies $c$,
which smooths per-frequency noise. $K_c$, by contrast, is the \emph{median} Pearson
correlation coefficient over individual per-$c$ curves; for slow-mixing Hamiltonians
with competing KAM structure, the per-$c$ variance is high and the per-$c$
log--log curves are non-monotonic, driving the median correlation negative even
though the mean $M_c$ grows with slope near one.
Yang--Mills shows $\alpha=-0.08$ (bounded $M_c$ at this run length);
H\'{e}non--Heiles ($\alpha=0.95$, $K=0.51$), double pendulum ($\alpha=1.31$,
$K=0.79$), and coupled Duffing ($\alpha=0.76$, $K=0.92$) show both measures
consistently positive.

\paragraph{Quasiperiodic length dependence.}
The quasiperiodic 2-torus achieves $K=0.961$ at $N=4{,}000$ steps
(Table~\ref{tab:systems}) but $K=-0.06$ and $\alpha=0.91$ at $N=12{,}000$ steps.
This confirms that the high $K$ is a finite-trajectory artefact: near-resonant $c$
values produce a visible $n^2$ transient that the log--log correlation records as
near-one at short runs; longer trajectories reveal saturation.
The long-range sub-slope $\alpha_{\mathrm{lo}}=0.37$ (partial saturation) versus
$\alpha_{\mathrm{lo}}=-0.81$ for the harmonic oscillator (unambiguous saturation)
shows that the sub-slope signature can flag periodic/quasiperiodic dynamics where
the full slope cannot. Neither $\alpha$ alone nor $K$ alone reliably separates
quasiperiodic from slow-mixing chaotic at finite trajectory lengths.

\begin{table}[ht]
\centering
\small
\caption{Growth-rate exponent $\alpha$ at $N=12{,}000$ steps. $\alpha$:
full-range log--log slope of $M_c(n)$; $\alpha_{\mathrm{lo}}$: long-range
($n\ge N_{\mathrm{dec}}/2$) sub-slope; $K$: 0--1 correlation from the same
12,000-step run (compare Table~\ref{tab:systems} for 4,000-step values; the
decimation stride cap increases from 20 to 60 with the longer trajectories,
which can lower $K$ for some dissipative flows without affecting $\alpha$).
Negative $\alpha_{\mathrm{lo}}$ signals $M_c$ saturation. Bold: Hamiltonian
systems with $\alpha\approx1$ yet $K\approx0$ (the $\alpha$--$K$ decoupling;
see text). Full results in \path{results/chaos_slope_analysis.json}.}
\label{tab:slope}
\begin{tabular}{llrrr}
\toprule
System & Family & $\alpha$ & $\alpha_{\mathrm{lo}}$ & $K$ \\
\midrule
\multicolumn{5}{l}{\textit{Dissipative flows}} \\
Lorenz 63              & Dissipative & $0.929$ & $+0.934$ & $+0.999$ \\
R\"{o}ssler            & Dissipative & $0.615$ & $-0.254$ & $+0.567$ \\
Duffing (forced)       & Dissipative & $0.916$ & $+0.995$ & $+0.276$ \\
Lorenz 96 ($N=5$)     & Dissipative & $0.884$ & $+0.090$ & $+0.940$ \\
Aizawa                 & Dissipative & $0.853$ & $+0.191$ & $+0.696$ \\
Sprott B               & Dissipative & $0.956$ & $+0.969$ & $+0.997$ \\
Halvorsen              & Dissipative & $1.018$ & $+1.331$ & $+0.957$ \\
Thomas                 & Dissipative & $1.098$ & $+1.132$ & $+0.995$ \\
Chen                   & Dissipative & $1.024$ & $+1.103$ & $+0.994$ \\
Ueda (forced)          & Dissipative & $0.971$ & $+1.297$ & $+0.295$ \\
Lorenz 96 ($N=10$)    & Dissipative & $0.902$ & $+0.978$ & $+0.998$ \\
\midrule
\multicolumn{5}{l}{\textit{Hamiltonian flows}} \\
Double pendulum                  & Hamiltonian & $1.312$ & $+1.296$ & $+0.792$ \\
H\'{e}non--Heiles                & Hamiltonian & $0.949$ & $+1.111$ & $+0.509$ \\
\textbf{CR3BP (Earth--Moon)}     & Hamiltonian & $\mathbf{0.847}$ & $+0.887$ & $\mathbf{-0.099}$ \\
Yang--Mills ($x^2y^2$)           & Hamiltonian & $-0.083$ & $+0.275$ & $+0.094$ \\
\textbf{Pullen--Edmonds}         & Hamiltonian & $\mathbf{0.694}$ & $-0.149$ & $\mathbf{-0.172}$ \\
\textbf{Quartic coupled}         & Hamiltonian & $\mathbf{0.888}$ & $+1.061$ & $\mathbf{-0.071}$ \\
\textbf{Spring pendulum}         & Hamiltonian & $\mathbf{0.972}$ & $+0.924$ & $\mathbf{-0.096}$ \\
Coupled Duffing (Ham.)           & Hamiltonian & $0.759$ & $+0.912$ & $+0.917$ \\
\midrule
\multicolumn{5}{l}{\textit{Maps}} \\
Logistic ($r=3.9$)    & Map   & $1.011$ & $+1.107$ & $+0.999$ \\
H\'{e}non map         & Map   & $1.020$ & $+1.008$ & $+1.000$ \\
Standard map          & Map   & $1.051$ & $+1.129$ & $+0.998$ \\
\midrule
\multicolumn{5}{l}{\textit{Delay system}} \\
Mackey--Glass         & Delay & $1.400$ & $+1.539$ & $+0.805$ \\
\midrule
\multicolumn{5}{l}{\textit{Non-chaotic controls}} \\
Harmonic oscillator   & Control & $0.543$ & $-0.806$ & $-0.111$ \\
Quasiperiodic 2-torus & Control & $0.906$ & $+0.373$ & $-0.063$ \\
Random walk           & Control & $0.008$ & $+0.162$ & $+0.140$ \\
\bottomrule
\end{tabular}
\end{table}

\section{Parameter sweeps and stress checks}

Parameter sweeps stress-test the two conclusions (Fig.~\ref{fig:mtsweep}). For
Lorenz63, $\nabla\!\cdot f=-\sigma-1-\beta$ is independent of $\rho$; the measured
$\Sigma\lambda$ remains pinned at $-13.667$ across $\rho\in[28,150]$ while the
largest Lyapunov exponent varies by a factor of over 33. Thus volume contraction
measures dissipative structure, not chaos strength. The 0--1 test, by contrast,
correctly flags non-chaotic regimes within a family: a Lorenz periodic window and
a standard-map KAM regime produce low $K$. This supports the corrected
interpretation: 0--1 is a chaos detector whose finite-time behavior depends on
mixing, while $\Sigma\lambda$ is the conservative/dissipative discriminator when
the vector field is available.

\begin{figure}[h]
\centering
\includegraphics[width=0.95\textwidth]{fig4_mixing_time_and_parameter_sweep}
\caption{Left: mixing-time law, with longer observable autocorrelation associated
with later $K$ crossover. Right: Lorenz volume contraction remains pinned to the
analytic trace while $\lambda_1$ changes substantially across the parameter sweep.}
\label{fig:mtsweep}
\end{figure}

\section{Discussion}
\label{sec:disc}

The practical message is simple. If one only has a scalar time series and wants to
ask whether the dynamics are chaotic, the 0--1 test remains useful. If one has a
model vector field and wants to distinguish conservative from dissipative chaos,
volume contraction is the theory-grounded diagnostic. The failure mode arises when
a finite-time scalar chaos detector is quietly promoted into a structural
classifier. Slow mixing then masquerades as conservative structure.

A striking secondary illustration of this failure: the quasiperiodic 2-torus
control (Table~\ref{tab:systems}) achieves $K=0.961$ at 4000 steps---higher than
five of the eight chaotic Hamiltonian systems. Quasiperiodic orbits produce
$M_c(n)\sim n^2$ (ballistic growth) near resonance, which the log--log form of
$K_c$ records as near-one. This is not chaos detection; it is a false positive
from a non-chaotic near-resonant system. Extended to 12,000 steps the same system
yields $K=-0.06$ (Section~\ref{sec:slope}), confirming that the high value is a
finite-trajectory artefact: the $n^2$ transient saturates before the longer
run ends. The result reinforces the recommendation to use $K$ as a \emph{chaos
detector}, not as a structural classifier.

When the model vector field is unavailable---purely observational scalar time
series---distinguishing conservative from dissipative chaos without computing
$\Sigma\lambda$ remains an open problem. The Green--Kubo framework (Section~\ref{sec:gk})
characterizes why the 0--1 statistic fails, but does not resolve it.

The itinerary result points in a different direction. It does not overcome the
Lyapunov wall for trajectories; it changes the prediction target. Local microstate
forecasts fail at long horizons, but coarse symbolic route grammar can remain
predictable on held-out data. That distinction is useful for scientific machine
learning systems because it separates impossible weather-like prediction from
bounded climate-like structure.

A notable pattern within the itinerary data: Hamiltonian systems achieve a higher
mean held-out transition accuracy ($0.841$, $n=8$) than dissipative flows ($0.746$,
$n=10$). This is the opposite of the naive expectation that strange attractors (with
their lower-dimensional structure) would be easier to learn. The likely explanation
is that several Hamiltonian systems in the benchmark are semi-ergodic: they orbit
near KAM tori rather than filling the energy surface uniformly, producing highly
predictable transition sequences. Four of the eight reach perfect held-out accuracy
($1.000$). This suggests that the grammar learnability correlates more strongly with
the regularity of the attractor's coarse geometry than with the sign of $\Sigma\lambda$.

An additional note on the random walk control: Brownian motion (the position
process) gives $K=0.045$ because the non-stationarity of the position observable
causes $M_c(n)$ to grow anomalously slowly under the log--log measure. The 0--1
test correctly identifies this as non-chaotic, but for a mechanically different
reason than a bounded regular orbit.

\section{Limitations and claim boundary}

Supported claims are limited to the 26-system benchmark and the deterministic
artifacts listed in Section~\ref{sec:repro}. The positive discriminator claim is empirical over 8
Hamiltonian systems and 11 dissipative finite-dimensional flows with defined vector
fields; Liouville's theorem supplies the theoretical reason for Hamiltonian zero
contraction, but this paper does not prove a new theorem. The cautionary 0--1 claim
is stronger in practice: the observed length dependence and Green--Kubo prediction
show why the statistic should not be used as a conservative/dissipative classifier.
The itinerary result is held-out symbolic transition prediction, not microstate
forecasting and not a claim about general trajectory prediction.

Retracted: an earlier version of this analysis incorrectly concluded that
conservative chaos has $K\approx0$; that reading is superseded here.
Not supported: using the 0--1 statistic as a conservative/dissipative classifier;
generalizing the benchmark to all chaotic systems without further validation;
claiming a new theorem; or claiming long-horizon chaotic trajectory prediction.
The manuscript intentionally preserves the self-disproved 0--1 result and the
corrected claim boundary rather than deleting the error.

\section{Reproducibility and data availability}
\label{sec:repro}

All reported numbers are read from deterministic, seeded artifacts. Data and figure-generation scripts are available at
\url{https://github.com/MonongahelaHellbender/chaos-universality-paper}. The primary evidence files are:
\begin{quote}
\small\raggedright
\path{results/chaos_universality_lab_best.json};
\path{results/chaos_mixing_time_analysis.json};
\path{results/chaos_green_kubo.json};
\path{results/chaos_parameter_sweep.json};
\path{results/chaos_itinerary_prediction.json}; and
\path{results/chaos_slope_analysis.json}.
\end{quote}
Figures are regenerated by \path{scripts/chaos_make_figures.py}.
The slope analysis (Table~\ref{tab:slope}) is regenerated by
\path{scripts/chaos_slope_analysis.py}.

\appendix

\section{Effective-sample interpretation of the 0--1 crossover}

Equation~\eqref{eq:mc} implies that $K$ reaches its asymptotic chaotic value only
when the trajectory is long relative to the observable's integrated
autocorrelation time. The artifacts estimate this time by a Sokal-style windowed
autocorrelation procedure \cite{sokal1997}. Across the cross-system ladder,
Spearman $r_s(\tau,n_{1/2})=0.58$ ($t=3.0$, $\mathrm{df}=17$, $p<0.01$; $n=19$
chaotic systems with measurable crossover; 4 Hamiltonian systems whose $K$ never
exceeded 0.5 within the audit range are excluded):
longer observable mixing time delays the first length where $K>0.5$.
Dissipative systems have median crossover near 3k steps,
while Hamiltonian systems have median crossover near 12k steps. The relevant
variable is therefore effective sample size after decorrelation, not conservation
versus dissipation.

\section{Reference table for the main claim}

\begin{center}
\small
\begin{tabular}{>{\raggedright\arraybackslash}p{0.25\linewidth}>{\raggedright\arraybackslash}p{0.32\linewidth}>{\raggedright\arraybackslash}p{0.34\linewidth}}
\toprule
Question & Correct diagnostic & Boundary \\
\midrule
Is the scalar series chaotic? & 0--1 $K$ & Requires sufficient effective length \\
Is the flow conservative or dissipative? & $\Sigma\lambda=\langle\nabla\!\cdot f\rangle$ & Requires vector field/model \\
Can the exact future state be predicted? & No long-horizon claim & Lyapunov wall remains \\
Can coarse route grammar be predicted? & Held-out itinerary model & Symbolic transitions only \\
\bottomrule
\end{tabular}
\end{center}

\begin{thebibliography}{99}

\bibitem{lorenz1963}
E.~N. Lorenz.
\newblock Deterministic nonperiodic flow.
\newblock \emph{Journal of the Atmospheric Sciences}, 20(2):130--141, 1963.

\bibitem{henon1964}
M.~H\'{e}non and C.~Heiles.
\newblock The applicability of the third integral of motion: Some numerical experiments.
\newblock \emph{Astronomical Journal}, 69:73--79, 1964.

\bibitem{mackey1977}
M.~C. Mackey and L.~Glass.
\newblock Oscillation and chaos in physiological control systems.
\newblock \emph{Science}, 197(4300):287--289, 1977.

\bibitem{strogatz1994}
S.~H. Strogatz.
\newblock \emph{Nonlinear Dynamics and Chaos}.
\newblock Addison-Wesley, 1994.

\bibitem{gottwald2004}
G.~A. Gottwald and I.~Melbourne.
\newblock A new test for chaos in deterministic systems.
\newblock \emph{Proceedings of the Royal Society A}, 460:603--611, 2004.

\bibitem{gottwald2009}
G.~A. Gottwald and I.~Melbourne.
\newblock On the implementation of the 0--1 test for chaos.
\newblock \emph{SIAM Journal on Applied Dynamical Systems}, 8(1):129--145, 2009.

\bibitem{falconer2007}
I.~Falconer, G.~A. Gottwald, I.~Melbourne, and K.~Wormnes.
\newblock Application of the 0--1 test for chaos to experimental data.
\newblock \emph{SIAM Journal on Applied Dynamical Systems}, 6(2):395--402, 2007.

\bibitem{benettin1980}
G.~Benettin, L.~Galgani, A.~Giorgilli, and J.-M. Strelcyn.
\newblock Lyapunov characteristic exponents for smooth dynamical systems and for Hamiltonian systems.
\newblock \emph{Meccanica}, 15:9--20, 1980.

\bibitem{kantz2004}
H.~Kantz and T.~Schreiber.
\newblock \emph{Nonlinear Time Series Analysis}.
\newblock Cambridge University Press, second edition, 2004.

\bibitem{green1954}
M.~S. Green.
\newblock Markoff random processes and the statistical mechanics of time-dependent phenomena.
\newblock \emph{Journal of Chemical Physics}, 22:398--413, 1954.

\bibitem{kubo1957}
R.~Kubo.
\newblock Statistical-mechanical theory of irreversible processes. I.
\newblock \emph{Journal of the Physical Society of Japan}, 12:570--586, 1957.

\bibitem{sokal1997}
A.~D. Sokal.
\newblock Monte Carlo methods in statistical mechanics: foundations and new algorithms.
\newblock In \emph{Functional Integration: Basics and Applications}, 1997.

\end{thebibliography}

\end{document}
